\chapter{55}



«Der da ist Onkel Moses, der Zuchtbock,» Mazi wies auf eine riesige,
bewegliche Wollkugel, die vor ihnen stehen blieb. Louis entdeckte
mittendrin zwei wache Augen. Sie fixierten ihn. Mazi sprach jedes Tier
mit Namen an. Louis konnte die Schafe beim besten Willen kaum
unterscheiden.
Schliesslich hatten sie alle aufgespürt. Sie waren wohlauf.
Mit ihrer zugewandten Art strahlten sie eine Sanftmut aus, die Louis
nicht erwartet hatte.

«Bei den trächtigen Mutterschafen musst du speziell gut hinschauen.
Toffee komm, komm her, Toffee,» mit den Worten näherte sich Mazi einem
Schaf mit straffem, beeindruckend dickem Bauch. Das Tier trottete ihm
bedächtig entgegen.

«Sollte es sich ungewöhnlich bewegen, Geräusche von sich geben, die dir
sonderbar erscheinen, oder misstrauisch auf dich reagieren, meldest du
dich unverzüglich bei Manuela. Bei Blut- oder Wasserverlust über die
Scheide, schlimmer noch bei Schleimabgang, musst du besonders schnell
reagieren. Bis jetzt erlebte ich einen einzigen solchen Notfall. Manuela
musste mit dem Heli hochgeflogen werden, um das Schaf zu holen. Es
konnte unten auf dem Hof verarztet werden.»

Louis wurde je länger je ungemütlicher. Mazis Erklärungen warfen ihn aus dem
Sattel. Jetzt, wo er sich mit seinen neuen Aufgaben gerade anfreundete,
kam noch ein gutes Stück Verantwortung oben drauf. Mazi blickte schräg
zu ihm hoch und schmunzelte.

«Die meiste Zeit geschieht nichts dergleichen. Ich bin jederzeit da,
wenn du dich unsicher fühlst,» er klopfte Louis auf die Schulter, «ruf
mich einfach an, okay?»

Mit zügigen Schritten verliess er Louis, drehte den Kopf zurück und
winkte ihm zu.

«Ich freue mich auf den Abend.»

\sectionbreak

Louis trug den Nachtisch nach draussen. Mazi hielt sich den Bauch.

«Das war ja gut. Ich weiss nicht, ob da noch Platz ist.»

Er verzog den Mund zu einer Grimasse und Louis lachte.

«Du scheinst Joh ganz schön begeistert zu haben. Er sagt, du bist
Musiker?»

Louis wischte sich über den Mund, setzte sich und griff nach seinem
Glas.

«Erst noch in Entwicklung, könnte man sagen. Meine eigenen Kreationen
sind noch nicht ausgereift. Mein Geld verdiene ich mit Coverversionen an
kleinen Auftritten und mit Strassenmusik.»

Ohne Pause holte Louis zur Gegenfrage aus. Weitere Lügen wollte er Mazi
ersparen.

«Was ist das für ein Song, den Joh zu seiner \emph{Alpmusik} auserkoren
hat? - wie mir Manuela sagte. Sie meinte, die afrikanische Musik hättest \emph{du} ihm
zugespielt.»

Mazi setzte sich gerade auf. Louis war, als wachse der gerade ein Stück.
Mit einem breiten Lächeln hellte sich sein Gesicht auf.

«Ja, \emph{«Manitoumani»}, meinst du, von Chedid - ein toller
Musiker,» damit zog er sein Handy aus der Hosentasche, «der schloss 
sich für einige Songs mit seinen afrikanischen Freunden zusammen. Ich
kenne die Drehorte, diese magische Umgebung. Die Einheimischen sind
grossartig. Ich fühl mich ihnen sehr verbunden. Hast du Lust
reinzuhören?»

Mazis Augen blitzten.
Louis nickte erwartungsvoll. Endlich würde er Johs «Alpsong» zu hören
bekommen.
Mazi reichte ihm die Kopfhörer über den Tisch, suchte nach dem Titel und
blickte mit schmal geöffnetem Mund in Louis’ Gesicht. Er drückte die
Playtaste.


\qrblock{Matthieu Chedid, Toumani Diabate\\ Sidiki Diabaté, Fatoumata Diawara\\ \emph{«Manitoumani»}}{media/QR-Codes/038_Manitoumani_-M-,_Toumani_Diabaté,_Sidiki_Diabat,_Fatoumata_Diawara_Spotify.png}

«Musikalisch wunderschön, gefällt mir sehr gut,» Louis legte die
Kopfhörer auf den Tisch, «spezieller Text, nicht?»

Mazi runzelte die Stirn. Er liess sich
Zeit.

«Ja, die Worte drücken Wahrheit aus, ohne hart zu werden. Die typisch
afrikanische Weise, unangenehme Situationen zu beschreiben. Meine
afrikanische Hälfte freut sich,» er stand auf und drehte sich unerwartet
um sich selbst, «der Text hat was Lyrisches und zusammen mit der Musik,» Mazi
schnalzte mit der Zunge und klatschte in die Hände, «einfach schön, wie
du sagtest.»

Louis war begeistert. Mazi war so spritzig anders. Er dachte 
an die endlosen Diskussionen der letzten Monate um die italienischen
Parlamentswahlen. Um Politik überhaupt. Eifrig und laut brachten Louis 
und seine Clique ihre klugen Lösungen vor. Regelmässig endeten die
Gespräche im Durcheinander. Sie fielen einander ins Wort und dann folgte
eine Runde Jammern ums Geld. Um die niedrigen Löhne und überhöhten
Preise. Und am Ende war alles genauso wie vorher. Mit Mazi fühlte sich Louis auffallend leichter.

\sectionbreak  

Die Sonne verschwand kurz nach acht. Mazi war im Halbdunkel abgestiegen.
Über der Umgebung lag Eintracht. Louis’ Körper fühlte sich angenehm müde
an.
Er räumte die Küche auf. Schliesslich blieb er vor dem kleinen
Südfenster stehen und schaute hinaus in die Nacht. Lange war es her,
seit er einen ganzen Tag in der Natur verbracht hatte. Als Kind hatte er
stunden-, wenn nicht tagelang, draussen gespielt. Er erinnerte sich an
die Nachmittage bei Grandpère zuhause, in diesem kleinen Wäldchen mit
all den Kindern der Siedlung. Wie er jeweils frisch durchlüftet und
erschöpft zurückgekehrt war. In Frieden mit sich. Wie gerade jetzt. Louis
hielt inne, stellte einen letzten, abgetrockneten Teller ins Regal und
stützte sich auf dem Tisch ab.
Die Menschen hier oben tickten auffallend anders. Ob er sich etwas
vormachte, sich täuschte? Es würde sich zeigen.
Mit den Gedanken legte er sich ins Bett, drehte sich auf die Seite und
klemmte die Decke zwischen seine Beine. Er atmete entspannt aus, als
plötzlich sein Handy klingelte. Überlaut, als wolle der Sound den
Frieden zerreissen. Er zuckte zusammen. Er hatte vergessen, die
Lautstärke zu dimmen. Louis starrte die Zahlen an, erinnerte und
entspannte sich. Joh.
